% page 302 of the facsimile (image p-301 of the scan)
% block 160.0 x 237.7 mm ; left margin 8.0 mm
\pgc{-12.700mm}{0.000mm}{
\l{\cel{3}294.}
\l{}
\l{\sou{konestablo}. (\sou{olim}). I. Oficiro di la reji di Francia. -}
\l{\cel{3}II. Chef-mayoro di la armei rejala. - DEFIS.}
\l{}
\l{\sou{konaxo}. (\sou{mat.}) Equaciono algebrala homogena inter x, y, z}
\l{\cel{3}e\cel{9}- x, y, z\cel{2}esante la koordinati homogena di}
\l{\cel{3}punto --\cel{7}la koordinati homogena di rekto.}
\l{}
\l{\sou{konfecionar}. (\sou{trans.}) Facar (vesti) granda-serio, segun}
\l{\cel{3}mezuri ne individuala, ma generala. - DF.}
\l{}
\l{\sou{konfekto}. Bonbono qua su dissolvas en la boko, e qua kon-}
\l{\cel{3}tenas liquoro o pasto sukroza e parfumoza. - DEIS.}
\l{}
\l{\sou{konferar}. (\sou{netrans., kun)}. Konversar (kun ulu) quan onu}
\l{\cel{3}konsultas pri temo konvencionita. - DEFIRS.}
\l{}
\l{\sou{konfervo.}\cel{2}(\sou{biol}.) Speco di algo kun filamenti tubatra.}
\l{\cel{3}- L.\cel{1}\sou{conferva}. - DEFISL.}
\l{}
\l{\sou{konfesar}. (\sou{trans., ad)}. I. Agnoskar kom da su ulo quon onu}
\l{\cel{3}agis, partikulare ulo etiko-mala. - II. (\sou{teol. katol.})}
\l{\cel{3}Deklarar volunte (sua peci) koram la tribunalo di peni-}
\l{\cel{3}tenco. - EFIS.}
\l{}
\l{\sou{konfesiono}. (\sou{religio kristana}). Deklaro koram publiko (+pu-}
\l{\cel{3}blico) pri sua acepto di la kredendaji di ta o ca eklezio.}
\l{\cel{3}- DEFIS.}
\l{}
\l{\textquotedbl{}\sou{konfetti}\textquotedbl{}. Drajei, buleti de gipsajo, dina rondaji de pa-}
\l{\cel{3}pero koloroza, quin onu interlansas dum karnavalo. - FI.}
\l{}
\l{\sou{konfidar}. (\sou{trans., ad ulu)}. Donar sekure (ulu od ulo) ad}
\l{\cel{3}ulu quan onu komisas por sorgar lu. - EFIS.}
\l{}
\l{\sou{konfidencar}. (\sou{trans,}\cel{1}ad).Konfidar (ad ulu) ulo ne-mate-}
\l{\cel{3}riala quon la konfidario devas ne divulgor - pri quo}
\l{\cel{3}ta konfidario devas tacar (sekretajo). - DEFIRS.}
\l{}
\l{\sou{konfirmar}. (\sou{trans.}) I. Igar (ulu) mem plu ferma, plu certa}
\l{\cel{3}(+certena) (pri). Igar plu certa, plu nedubitebla, per pru-}
\l{\cel{3}vo. - II. Adjuntar a la vigoro quan (lu) ja posedas, afir-}
\l{\cel{3}mon, pruvon, e c. - III. (\sou{religio}) Donar ad (ulu) ta ek la}
\l{\cel{3}sep sakramenti di la eklezio katolika, qua havas kom skopo}
\l{\cel{3}igar la Santa Spirito decensar adsur la kristano por igar}
\l{\cel{3}plu ferma la gracio de bapto. - IV. (\sou{komerco})\cel{2}Repetar(la}
\l{\cel{3}texto di sua letro, telegramo, e c.) por ke la destinario}
\l{\cel{3}esez mem plu +certena (certa) pri lo afirmita. - V. (\sou{pri}\cel{1}in-}
\l{\cel{3}\sou{formo qua ne esis absolute certa}). Afirmar ke ta informo}
\l{\cel{3}esas nun certa, nedubitebla. - DEFIRS.}
\l{}
\l{\sou{konfiskar}. (\sou{trans.}) I. Atribuar a fisko, per dekreto, per}
\l{\cel{3}lego (to quon proprietas ulu). II. Embargar ulo quan}
\l{\cel{3}interdiktas la reguli di la skolo, di la liceo, e c.-DEFIRS.}
}
